\section{Discussion}
\label{sec:discussion}

We have described how the error of a case-count estimate of the reproduction number depends on case volume when transmission varies between periods, offspring numbers are overdispersed and only a proportion of infections is reported. For the period-specific value, error falls without limit as reported cases accumulate. For the mean across periods and for one-step prediction, error per period approaches the environmental variance, and the balance scale $\mx=A/v_R$ marks the case volume beyond which each doubling of cases removes at most a quarter of the remaining variance. In weekly influenza hospital admissions over four seasons, an affine model gave both apparent statistics positive intercepts, although an intercept-free power law fitted the observed range about as well. Conditional on a mean generation interval of 3.6 days and on independent weekly shocks, a two-step inversion through a data-matched simulation mapped both intercepts to a latent environmental variance of about 0.04. A national-week component accounted for about 40\% to 50\% of the intercept, and facility completeness under mandatory reporting contributed less than equivalence margins of 17\% for $R_A$ and 12\% for $G$, set at a reference reporting volatility of 0.10.

The theory carries a standard result in population ecology over to surveillance data, and its general form is not new. Ecological models split the variance of population growth into a demographic term that falls with population size and an environmental term that does not \citep{engen1998demographic, lande2003stochastic}. The present model retains this partition and adds four elements that matter for surveillance. Transmission is a branching process with negative binomial offspring, so that superspreading enlarges the demographic term \citep{lloydsmith2005superspreading}. Binomial reporting adds a further term that falls with case volume, and reporting volatility adds one that does not (Proposition~\ref{prop:report-noise}). The quantity of interest is an estimate averaged over periods, so the balance scale also sets the optimal division of a budget between cases per period and number of periods (Proposition~\ref{prop:allocation}). That rule has the form of optimal allocation in two-stage sampling \citep{cochran1977sampling} and of the optimal cluster size in cluster-randomised trials \citep{raudenbush1997statistical}, with the balance scale in place of the ratio of within-unit to between-unit variance. The contribution of the present model is therefore the mapping of these terms onto the reproduction number, together with the reporting term and the calibration that links apparent weekly statistics to the latent variance, and not the variance partition itself.

The results complement the information-theoretic analyses of Parag and colleagues \citep{parag2021improved, parag2022quantifying}. Their metric for a noisy epidemic curve already contains time-varying reporting fractions, which enter through the geometric mean of the reporting probabilities, and a negative binomial extension is treated in their supplementary material \citep{parag2022quantifying}. Their target is the sequence of period-specific values, and in a Poisson renewal model the error variance of each estimate is $\Rbar/m$, so that reporting acts by rescaling the effective case count. Here the target is the mean across periods, and variation shared by all chains of transmission in a period contributes variance that is independent of $m$. The two descriptions agree at small counts and separate near the balance scale. Under the affine fit on the apparent scale of $R_A$, the fitted variance was twice the Poisson variance at 22 admissions per week and ten times at 261, and on the latent scale the corresponding values were 28 and 262. Both descriptions are therefore needed. The Poisson analysis ranks data streams by the information they carry about current transmission, and the variance partition here indicates when further cases cease to improve estimates of typical transmission.

The two-step inversion connects the apparent statistics with the latent model. The simulation first established that neither statistic estimates $v_R/\Rbar^2$ directly; the intercept of $G$ was 2.7 times this quantity and that of $R_A$ about half of it, for reasons that follow from how each statistic aggregates growth over a week. After calibration, the two statistics implied latent variances of 0.043 and 0.040 (ratio 0.93; 95\% CI 0.78 to 1.18). This consistency is expected if both statistics respond to the same latent variation and the simulation captures how each transforms it. As a test of the calibration it has limited power, since across the simulated settings it failed only when a synchronous shock to reported counts replaced the shock to transmission, in which setting the observed lag-1 autocorrelation of $G$ lay above 99\% of simulated panels. The simulated panels also reproduced the intraclass correlation and the state-level lag-1 autocorrelation of the data; the national-week share was set from the data and is not a test. Within the stated uncertainties, the calibration is therefore consistent with the empirical intercepts, under three conditions.

The first condition is the generation interval. The 3.6 days assumed by the statistics is the mean household serial interval of \citet{cowling2009estimation}, used as a proxy for the generation interval. The interval of 2.9 to 4.3 days is the 95\% confidence interval of that mean, not the range of plausible true intervals. Direct estimates of the mean generation interval, 2.8 to 3.2 days \citep{chan2025estimating}, and mean serial intervals of 2.2 and 2.8 days for influenza A subtypes \citep{vink2014serial} lie in the lower part of this interval or below it. Across true mean intervals of 2.2 to 4.3 days the implied variance ranged from 0.013 to 0.063, and at 2.9 days it was 0.026. The inversion therefore identifies $v_R$ conditionally on the interval, and its primary interval does not include uncertainty about it. If the true interval is shorter than 3.6 days, the latent variance is smaller and the latent balance scale larger, about 47 admissions per generation at 2.9 days. The second condition is independence of weekly shocks. The pooled state-level autocorrelation lay near the centre of the simulated distribution under independence and below 97\% of simulated panels with a first-order coefficient of 0.25 for $R_A$. A coefficient of 0.5 would have lowered the implied variance to about 0.026. This check concerns autocorrelation within states; the series of national-week means was not tested, and the simulated autocorrelation acted on the national and state shocks together. The third condition concerns the cross-sectional slope. The calibrated theoretical $b/a$ lies below both empirical point estimates, which fall at the 91st and 93rd percentiles of the simulated sampling distribution. With a variance-to-mean ratio of 1.5 for weekly admissions, the simulated $b/a$ rose to 68.0 for $R_A$ and 40.2 for $G$, above both empirical values, while the implied variance stayed at 0.039 and 0.044. A true generation interval of 2.9 days also reproduced the empirical $b/a$, with an implied variance of 0.026. The excess can therefore come from overdispersion smaller than a variance-to-mean ratio of 1.5, from a shorter generation interval, or from spatial aggregation (below), and these explanations imply different latent variances. The state-level autocorrelation does not favour the shorter interval, but it cannot exclude it, and the data cannot exclude a mixture of transmission and reporting shocks (Proposition~\ref{prop:report-noise}).

The fit across states also confounds case volume with spatial scale. States with more admissions are, in general, larger and more populous, and their weekly counts aggregate more local epidemics. If local environmental variation is not fully synchronised within a state, Eq.~(\ref{eq:pooling}) implies that the environmental variance of a state-level statistic falls with the number of local epidemics it contains. The relative loss would then decline with state size even if the shared term within each local epidemic were constant, and part of the slope $b$ would reflect aggregation and not sampling. Under this explanation, $b/a$ overstates the balance scale within a single reporting unit, and the design rule below applies to state-level aggregates of the size observed here. Eq.~(\ref{eq:affine}) assumes a common intercept for all states, and the data do not separate the two explanations.

Under the affine model, the empirical results are best read as a diminishing-returns scale for case volume within a period. With an apparent balance ratio of 37.4 admissions per week on the scale of $R_A$, a state that already reports at this ratio obtains half of the precision that unlimited case volume could provide, and 90\% requires nine times as many admissions. Slightly more than half of the state-weeks in the analysis windows were above the apparent ratio; the latent scale implied by the calibration is higher, so fewer state-weeks lie above it. The limit applies to case volume per period. The error of an estimate pooled over $T$ independent periods falls as $v_R/T$, estimators that shrink towards a prior mean can go below the per-period value (Appendix~\ref{app:shrinkage}), and a power law without an intercept described the observed range of scales about as well as the affine form. The balance scale is therefore a statement about where returns diminish, and it does not imply that error stops falling with more data of every kind. This interpretation is also robust to the source of the intercept. Whatever mixture of transmission, testing and statistic construction it contains, a component that is shared within a period and does not fall with case volume limits the value of further cases per period and raises the value of further periods.

The multiscale decomposition locates this component. About 41\% of the intercept of $R_A$ and 51\% of that of $G$ was common to all states in a calendar week, and resampling weeks moved the estimated intercept 1.5 to 1.9 times as much as resampling states. A national-week component cannot be reduced by scaling up surveillance in a single state, because it enters every state's loss in that week. Pooling states removes the state-specific part, and in these data pooled variance approached the national-week component, about a quarter of single-state variance, with 25 to 33 states. This share is consistent with national climatic and seasonal transmission influences \citep{cauchemez2008estimating, shaman2010absolute}, and with national changes in testing or admission practice. It is also consistent with a common trend within seasons, such as a slowing of growth as the season progresses, which is shared by all states without being week-to-week environmental variation; the present decomposition cannot separate the two. Two-way inference is required for the same reason; intervals that resample states only were about two to five times too narrow for the intercept \citep{owen2007pigeonhole, cameron2011robust}.

The contrast between reporting regimes points towards a benefit of high facility participation but does not establish it. Under mandatory reporting, a median of about 93\% of hospitals reported. Adjustment for reporting completeness changed the intercept of $R_A$ by 2\%, with a one-sided upper bound of 5.3\%, and the change was equivalent to zero within a margin of 17\%. Most state-weeks in the windows had nearly complete participation, so this bound concerns small departures from full participation, and at a reference volatility of 0.05 equivalence was not shown. During the voluntary period of 2024, when the median share fell to about 46\%, the interquartile range of log growth was 1.58 times that of the same months in 2023. The interval for this ratio, 0.95 to 2.23, includes one, and the comparison rests on a single off-season. Several checks address the main confounders. The placebo comparison of 2025 with 2023 gave 1.01, matching on reporting scale and on calendar month preserved a ratio above one, and COVID-19 and RSV admissions showed ratios above one in the same off-season, which is therefore not an independent replication. These descriptive patterns are compatible with higher empirical volatility during periods of low and variable facility participation, though they do not identify a causal effect of the reporting mandate. The adjustment measures facility participation only. Testing intensity, coding and care seeking within reporting hospitals are not captured. The laboratory proxy shared 23\% of the intercept of $R_A$, but it does not discriminate between transmission and ascertainment, because laboratory growth carries its own noise and attenuates the shared fraction towards zero.

These results suggest a design rule for routine surveillance of the mean reproduction number. Where weekly admissions per jurisdiction already exceed the price-adjusted threshold $m^*$, additional resources reduce error more when spent on further periods and on further jurisdictions whose variation is not synchronised (Eq.~(\ref{eq:pooling})) than on further cases per period. The value of complete reporting depends on the reporting volatility a system actually has. In the hypothetical example of Section~\ref{sec:res-theory}, with a relative volatility of $\tau=0.15$, standardising reporting reduced error more than a tenfold increase in case volume per period. The completeness bound observed under mandatory reporting allows $\tau$ up to about 0.056 by the scaling against the reference volatility, or about 0.040 on the latent scale (Section~\ref{sec:meth-reporting}). At $\tau=0.040$ a tenfold increase in cases reduced the mean squared error from $15.2\times10^{-3}$ to $11.4\times10^{-3}$, against $14.6\times10^{-3}$ for standardised reporting, and at $\tau=0.056$ from $15.7\times10^{-3}$ to $11.9\times10^{-3}$, against $14.6\times10^{-3}$. Under mandatory reporting, case volume therefore remained the more effective lever in this example. The value of complete reporting relative to case volume grows with the volatility that incomplete participation introduces, which the voluntary period suggests may be larger but did not establish. Seasons are the limiting unit for statements about variation between seasons. With four seasons, no unbiased estimator of a between-season variance can have a relative standard error below 0.82 (Corollary~\ref{cor:season}), and the intercept ratios for alternative denominator weights did not discriminate between mechanisms, because epidemic growth with reporting noise alone already produced $q_1<1<q_2$.

The scope of these conclusions is set by the data, the model form and the calibration. The empirical results rest on one surveillance system, four influenza seasons and rule-defined early-growth windows. The window rule was revised once and the renewal-type ratio became the primary statistic during the analysis; Table~\ref{tab:s9} lists each specification with its role and date, and the headline estimates lie within each other's intervals. The positive intercept and every quantity derived from it are conditional on the affine form, which the data did not prefer over a power law. The intercept is fitted to state-level aggregates and pools all seasons around one mean, so between-season differences in mean growth and spatial aggregation both enter it. The inversion is conditional on the generation interval, for which a household serial interval served as proxy, and on the simulated delay, offspring and independence structure. The statistic $R_A$ is a transformation of weekly growth, not a posterior estimate of $R_t$ \citep{cori2013new}. The propositions assume environmental variation that is independent between periods and a fixed number of primary cases; the renewal formulation of Section~\ref{sec:meth-loss} shows that uncertainty in infection pressure adds a further term that does not fall with case volume. The equivalence margins use the sample mean of each statistic, the null ranges for the intercept ratios in rule windows come from the calendar-window simulation, and the bin error bars in Fig.~\ref{fig:floor}B ignore clustering. Fits with state effects or stratified by state size would separate case volume from spatial aggregation, and an inversion that integrates over uncertainty in the generation interval would remove the main conditioning. Simulations with autocorrelation in the national shock alone and with the window rule applied to simulated data would test the independence assumption more directly. Renewal estimates with an explicit admission delay, hierarchical models with week effects and independent measures of testing, and a second surveillance system with mandatory reporting and more seasons would extend the empirical base.
